% Options for packages loaded elsewhere
\PassOptionsToPackage{unicode}{hyperref}
\PassOptionsToPackage{hyphens}{url}
\documentclass[
  ignorenonframetext,
]{beamer}
\newif\ifbibliography
\usepackage{pgfpages}
\setbeamertemplate{caption}[numbered]
\setbeamertemplate{caption label separator}{: }
\setbeamercolor{caption name}{fg=normal text.fg}
\beamertemplatenavigationsymbolsempty
% remove section numbering
\setbeamertemplate{part page}{
  \centering
  \begin{beamercolorbox}[sep=16pt,center]{part title}
    \usebeamerfont{part title}\insertpart\par
  \end{beamercolorbox}
}
\setbeamertemplate{section page}{
  \centering
  \begin{beamercolorbox}[sep=12pt,center]{section title}
    \usebeamerfont{section title}\insertsection\par
  \end{beamercolorbox}
}
\setbeamertemplate{subsection page}{
  \centering
  \begin{beamercolorbox}[sep=8pt,center]{subsection title}
    \usebeamerfont{subsection title}\insertsubsection\par
  \end{beamercolorbox}
}
% Prevent slide breaks in the middle of a paragraph
\widowpenalties 1 10000
\raggedbottom
\AtBeginPart{
  \frame{\partpage}
}
\AtBeginSection{
  \ifbibliography
  \else
    \frame{\sectionpage}
  \fi
}
\AtBeginSubsection{
  \frame{\subsectionpage}
}
\usepackage{iftex}
\ifPDFTeX
  \usepackage[T1]{fontenc}
  \usepackage[utf8]{inputenc}
  \usepackage{textcomp} % provide euro and other symbols
\else % if luatex or xetex
  \usepackage{unicode-math} % this also loads fontspec
  \defaultfontfeatures{Scale=MatchLowercase}
  \defaultfontfeatures[\rmfamily]{Ligatures=TeX,Scale=1}
\fi
\usepackage{lmodern}
\ifPDFTeX\else
  % xetex/luatex font selection
\fi
% Use upquote if available, for straight quotes in verbatim environments
\IfFileExists{upquote.sty}{\usepackage{upquote}}{}
\IfFileExists{microtype.sty}{% use microtype if available
  \usepackage[]{microtype}
  \UseMicrotypeSet[protrusion]{basicmath} % disable protrusion for tt fonts
}{}
\makeatletter
\@ifundefined{KOMAClassName}{% if non-KOMA class
  \IfFileExists{parskip.sty}{%
    \usepackage{parskip}
  }{% else
    \setlength{\parindent}{0pt}
    \setlength{\parskip}{6pt plus 2pt minus 1pt}}
}{% if KOMA class
  \KOMAoptions{parskip=half}}
\makeatother
\setlength{\emergencystretch}{3em} % prevent overfull lines
\providecommand{\tightlist}{%
  \setlength{\itemsep}{0pt}\setlength{\parskip}{0pt}}
% Auto-generated by infrastructure/rendering/_pdf_latex_helpers.py::extract_math_font_preamble.
% Minimal header passed to Pandoc via -H so Beamer slides pick up the same math font
% as the combined-PDF path. Adjust the manuscript's preamble.md to change the font globally.
\usepackage{unicode-math}
\setmathfont{latinmodern-math.otf}

% Slide layout defaults for warning-clean scientific decks.
\usepackage{etoolbox}
\IfFileExists{xurl.sty}{\usepackage{xurl}}{}
\IfFileExists{seqsplit.sty}{\usepackage{seqsplit}}{\newcommand{\seqsplit}[1]{#1}}
\protected\def\breakseq#1{\seqsplit{#1}}
\protected\def\breaktt#1{\begingroup\ttfamily\seqsplit{#1}\endgroup}
\setlength{\emergencystretch}{6em}
\tolerance=5000
\hbadness=10000
\hfuzz=1pt
\setlength{\tabcolsep}{2pt}
\AtBeginEnvironment{longtable}{\tiny\renewcommand{\arraystretch}{0.86}\setlength{\tabcolsep}{1pt}}
\AtBeginEnvironment{tabular}{\tiny\renewcommand{\arraystretch}{0.86}\setlength{\tabcolsep}{1pt}}

% Natbib and cross-reference fallbacks — slides don't load natbib
% or cleveref, but combined-PDF manuscript prose may emit these
% commands. The fallback renders citations as a bracketed key list
% and cross-references as detokenized labels so slides don't fail on
% undefined control sequences or raw underscores. \providecommand is
% a no-op if packages load later.
\providecommand{\citep}[1]{[#1]}
\providecommand{\citet}[1]{#1}
\providecommand{\citealp}[1]{#1}
\providecommand{\citeauthor}[1]{#1}
\providecommand{\citeyear}[1]{#1}
\providecommand{\cref}[1]{\texttt{\detokenize{#1}}}
\providecommand{\Cref}[1]{\texttt{\detokenize{#1}}}
\usepackage{bookmark}
\IfFileExists{xurl.sty}{\usepackage{xurl}}{} % add URL line breaks if available
\urlstyle{same}
% fallback for those not using the hyperref driver hyperxmp:
\makeatletter
\@ifundefined{xmpquote}{\newcommand{\xmpquote}[1]{#1}}{}
\makeatother
\hypersetup{
  hidelinks,
  pdfcreator={LaTeX via pandoc}}

\author{\texorpdfstring{}{}}
\date{}

\begin{document}

\section{Methodology}\label{sec:methodology}

\begin{frame}[fragile,allowframebreaks]{Methodology}
This section maps each step of the exploratory analysis to the tested
function that implements it. Every function lives in \texttt{src/eda/},
takes a \breaktt{pandas.DataFrame} in, and returns data out --- no
plotting, no file I/O --- so the notebook, the analysis script, and this
manuscript all reach the same code.
\end{frame}

\begin{frame}[fragile,allowframebreaks]{Dataset loading and schema}
\protect\phantomsection\label{dataset-loading-and-schema}
\breaktt{src/eda/dataset.py::load\_dataset()} reads the shipped CSV and
coerces the numeric columns with
\breaktt{pandas.to\_numeric(errors="coerce")}. Coercion is deliberate: a
blank or non-numeric cell becomes \texttt{NaN} rather than raising or
being silently dropped, so missingness is \emph{surfaced} and handled in
an explicit later step. The column roles are declared once in a frozen
\texttt{DatasetSchema} (an identifier column, a categorical group
column, and three numeric features), which downstream functions consult
instead of re-sniffing dtypes.
\end{frame}

\begin{frame}[fragile,allowframebreaks]{Cleaning: explicit, reported row
removal}
\protect\phantomsection\label{cleaning-explicit-reported-row-removal}
\breaktt{src/eda/cleaning.py::clean\_dataset()} drops any row missing a
numeric feature and returns a \texttt{CleaningReport} recording how many
rows entered, how many remained, and how many were dropped. This
exemplar makes a deliberate choice --- listwise deletion with a visible
count --- rather than imputation, because a first EDA pass should
\emph{see} the missingness, not paper over it. A separate
\breaktt{normalize\_numeric()} z-score standardizes the numeric columns
(mean 0, sample standard deviation 1), with a guard that maps a constant
or single-row column to zeros instead of producing
\texttt{inf}/\texttt{NaN}.
\end{frame}

\begin{frame}[fragile,allowframebreaks]{Descriptive statistics}
\protect\phantomsection\label{descriptive-statistics}
\breaktt{src/eda/statistics.py::summary\_statistics()} returns one
\texttt{ColumnSummary} per numeric column --- count of non-missing
observations, mean, sample standard deviation, minimum, median, and
maximum --- computed with real pandas aggregation.
\texttt{group\_means()} returns the mean of each numeric feature grouped
by the categorical column, sorted by group name for deterministic
output.
\end{frame}

\begin{frame}[fragile,allowframebreaks]{Correlation structure}
\protect\phantomsection\label{correlation-structure}
\breaktt{src/eda/correlation.py::correlation\_matrix()} returns the
Pearson correlation matrix of the numeric columns \breakseq{[@tukey1977eda]}.
The companion \breaktt{strongest\_pairs(matrix,\ top\_n)} ranks the
distinct off-diagonal feature pairs by absolute correlation while
preserving sign --- usually the single most useful artifact of a first
pass, because it points directly at the relationships worth
investigating. Each unordered pair appears exactly once.
\end{frame}

\begin{frame}[fragile,allowframebreaks]{Figure-data preparers}
\protect\phantomsection\label{figure-data-preparers}
Plotting is kept out of the library entirely.
\breaktt{src/eda/figures.py} returns \emph{plot-ready data structures}
--- bin counts and edges for a histogram, a square value grid plus
labels for a correlation heatmap, and sorted category counts for a bar
chart --- as frozen dataclasses of plain numbers. The thin analysis
script (\breaktt{scripts/eda\_analysis.py}) and notebook cells consume
these structures and call matplotlib. This keeps the library importable
on a headless machine and makes every preparer testable without a
display backend.
\end{frame}

\begin{frame}[fragile,allowframebreaks]{Zero-mock testing methodology}
\protect\phantomsection\label{zero-mock-testing-methodology}
The project is governed by a strict zero-mock policy, evaluated by
running
\breaktt{uv\ run\ pytest\ projects/templates/template\_eda\_notebook/tests}
during the build.

\begin{enumerate}
\tightlist
\item
  \textbf{Library tests} exercise every public function against the
  shipped CSV or a tiny hand-built frame with values chosen so the
  expected statistic is exact (e.g.~\breaktt{weight\ =\ 2\ *\ height}
  gives a correlation of exactly \texttt{+1.0}). No
  \texttt{unittest.mock}, no \texttt{MagicMock}, no \texttt{@patch}.
\item
  \textbf{Script test} runs \texttt{run\_eda()} against a temporary
  output root and asserts that real PNG figures and a real summary CSV
  are written.
\item
  \textbf{Notebook test} parses the real \texttt{.ipynb} and asserts it
  is valid nbformat, that every name imported \texttt{from\ src} exists
  in the library's public surface, and that no cell defines its own
  \texttt{def}/\texttt{class}.
\item
  \textbf{Coverage gate}: CI enforces a ≥90\% statement-coverage gate on
  \breaktt{projects/templates/template\_eda\_notebook/src/}; the live
  figure is tracked in
  \href{../../../../docs/_generated/COUNTS.md}{\breaktt{docs/\_generated/COUNTS.md}}.
\end{enumerate}
\end{frame}

\begin{frame}[fragile,allowframebreaks]{Figure generation contract}
\protect\phantomsection\label{figure-generation-contract}
Each figure in \texttt{03\_results.md} maps to a figure-data preparer in
\breaktt{src/eda/figures.py}: \texttt{histogram\_data} → height
histogram, \breaktt{correlation\_heatmap\_data} → feature-correlation
heatmap, and \breaktt{group\_count\_data} → per-group row counts.
Captions name the preparer and the key parameters (bin count, value
range) so reviewers can navigate from the PDF to the code without
inferring hidden defaults.
\end{frame}

\end{document}
